\subsection{Alternatif Mekanisme Data Integrity dan Non-Repudiation}

Terdapat beberapa pendekatan kriptografi yang dapat dipertimbangkan untuk menjamin integritas data dan \textit{non-repudiation} pada entri PGHD dalam ekosistem DecMed.
Parameter-parameter berikut digunakan sebagai dasar perbandingan mekanisme, mengikuti pola evaluasi yang diterapkan pada komponen lain di DecMed:

\begin{enumerate}
    \item \textbf{Jaminan \textit{non-repudiation}.} 
    Kemampuan mekanisme untuk membuktikan secara kriptografis bahwa data dibuat oleh entitas tertentu, sehingga entitas tersebut tidak dapat menyangkal kepengarangan data tersebut (\textit{non-repudiation}).
    \item \textbf{Jaminan integritas data.} 
    Kemampuan mekanisme untuk mendeteksi setiap modifikasi tidak sah terhadap data setelah data ditandatangani \textit{data integrity}.
    \item \textbf{Dapat diverifikasi tanpa kunci rahasia pasien.} 
    Kemampuan mekanisme untuk diverifikasi oleh pihak ketiga mana pun yang hanya memiliki kunci publik, tanpa memerlukan akses terhadap kunci privat atau kunci rahasia bersama.
\end{enumerate}

Tabel~\ref{tab:perbandingan-mekanisme-security} menyajikan perbandingan ketiga mekanisme berdasarkan parameter evaluasi yang telah ditetapkan.

\begin{table}[htbp]
    \centering
    \caption{Perbandingan mekanisme integritas data dan \textit{non-repudiation}}\label{tab:perbandingan-mekanisme-security}
    \renewcommand{\arraystretch}{1.25}
    \rowcolors{2}{white}{white}
    \begin{tabularx}{\textwidth}{>{\raggedright\arraybackslash}X
        >{\centering\arraybackslash}p{0.15\textwidth}
        >{\centering\arraybackslash}p{0.15\textwidth}
        >{\centering\arraybackslash}p{0.15\textwidth}}
        \rowcolor[HTML]{DDEBF7}
        \textbf{Mekanisme} & \textbf{P1} & \textbf{P2} & \textbf{P3} \\
        \midrule
        Hash saja (SHA-256) & Ya & Tidak & Tidak \\
        \midrule
        MAC / HMAC & Ya & Tidak & Tidak \\
        \midrule
        Digital Signature & Ya & Ya & Ya \\
        \bottomrule
    \end{tabularx}
    \bigskip

    \noindent\footnotesize
    \textbf{Keterangan:}
    P1: Jaminan integritas data;
    P2: Jaminan \textit{non-repudiation};
    P3: Dapat diverifikasi tanpa kunci rahasia.
\end{table}

Berdasarkan perbandingan pada Tabel~\ref{tab:perbandingan-mekanisme-security}, mekanisme \textit{digital signature} memakai pasangan kunci publik-privat (\textit{keypair}) pasien dipilih sebagai mekanisme yang akan diimplementasikan. %chktex 8
Pilihan ini didasarkan pada beberapa pertimbangan.

Pertama, hash kriptografis yang berdiri sendiri tidak dapat memenuhi kebutuhan \textit{non-repudiation} karena tidak mengikat identitas pengirim pada data.
Penyerang yang memodifikasi data dapat pula mengganti nilai hash secara bersamaan \parencite{banerjee2025blockchain}. 
Kedua, mekanisme MAC/HMAC bergantung pada kunci simetris yang dimiliki bersama oleh kedua pihak, sehingga tidak dapat memberikan jaminan \textit{non-repudiation}: mana pun dari kedua pihak dapat menghasilkan MAC yang valid \parencite{NIST_glossary_nonrepudiation}.

\textit{Digital signature} berbasis \textit{keypair} asimetris pasien merupakan pilihan yang paling sesuai karena memenuhi seluruh parameter: memberikan jaminan integritas dan \textit{non-repudiation} sekaligus, tidak bergantung pada infrastruktur terpusat apapun, dan dapat diverifikasi tanpa memerlukan \textit{private key} pasien atau key yang sama yang dipegang oleh pasien dan sistem DecMed \parencite{NIST_FIPS186-5, lei2023analysis, banerjee2025blockchain}. 
Proses penandatanganan dilakukan di sisi klien pasien menggunakan kunci privat, sedangkan verifikasi dilakukan oleh smart contract DecMed menggunakan kunci publik pasien yang telah terdaftar on-chain pada saat registrasi akun.